\section{金属の強化機構}

\subsection{令和2年度 問8}
\label{q:r02-reg-applied-08}

\noindent\textbf{原典：} sources/R02/17\_applied\_science.pdf、PDF物理3頁。

\textbf{公式正答：} 5（single）。

\paragraph{問題の要旨} 金属の強化機構を識別し、焼入れに関する変態名の誤りを見抜く。

\Needspace{4\baselineskip}\paragraph{確認した条件・記述（独自要約）}\mbox{}\par

\noindent\textbf{原典：} 指定頁を実画像で再確認。\par

\paragraph{問われる核心} 材料物性

\paragraph{必要知識} 転位、加工硬化、固溶・析出・粒界強化、マルテンサイト変態

\paragraph{考え方・正答を選ぶ理由（査読待ち初稿）} 鋼の焼入れ硬化で中心となるのはマルテンサイト変態である。⑤の「バウシンガー変態」は原典の誤答として保存し、解説中で正しい用語と区別する。したがって公式正答は⑤。

\Needspace{4\baselineskip}\paragraph{選択肢の検討（査読待ち初稿）}\mbox{}\par

\noindent\textbf{1：記述と範囲を確認} 加工硬化。塑性変形に伴う転位等により変形抵抗が増す。\par

\noindent\textbf{2：記述と範囲を確認} 固溶強化。溶質原子が格子・転位へ影響し強化する。\par

\noindent\textbf{3：記述と範囲を確認} 析出強化。固溶限の温度依存を利用した析出が転位運動を妨げる。\par

\noindent\textbf{4：記述と範囲を確認} 結晶粒微細化。粒界が転位の移動を妨げ通常降伏強さが増す。極端な粒径や温度まで無限定にしない。\par

\noindent\textbf{5：不適切・独立選択} 焼入れはバウシンガー変態。焼入れ硬化はマルテンサイト変態。Bauschingerは逆方向の降伏挙動に関する効果で焼入れ変態名ではない。\par

\paragraph{再利用する考え方} 転位、加工硬化、固溶・析出・粒界強化、マルテンサイト変態

\paragraph{共通理論} \hyperref[theory:material:strengthening]{第\ref*{theory:material:strengthening}節}

\par 難易度 2（校準初稿）、優先度未評価。
